\section{Mikołaj Łobaczewski i jego płaszczyzna} % lang-pl
\section{Nikolaj Lobachevskij e il suo piano} % lang-it
\label{sekcja_lobaczewski}

Opowiemy w tym rozdziale o geometrii Łobaczewskiego, nazwanej później przez Kleina geometrią hiperboliczną\footnote{Felix Klein w 1871 roku zaproponuje podział na geometrię eliptyczną (sferyczną), paraboliczną (Euklidesa) i hiperboliczną (Łobaczewskiego)}. % lang-pl
Mikołaj Iwanowicz Łobaczewski swoje idee przedstawi 23 lutego 1826 roku, a wydrukuje je w latach 1829--1830 w publikacjach uniwersytetu kazańskiego, oddzielony od świata naukowego, jak zauważy Eves, barierami odległości i języka \cite[s. 287]{eves1_1972}. % lang-pl
Wersję niemiecką, \emph{Geometrische Untersuchungen}, opublikuje w 1840 roku. % lang-pl
William Kingdon Clifford nazwie Łobaczewskiego ,,Kopernikiem geometrii''. % lang-pl
In questo capitolo parleremo della geometria di Lobachevskij, chiamata poi da Klein geometria iperbolica\footnote{Felix Klein nel 1871 proporrà la suddivisione in geometria ellittica (sferica), parabolica (di Euclide) e iperbolica (di Lobachevskij).}. % lang-it
Nikolaj Ivanovič Lobachevskij presenterà le sue idee il 23 febbraio 1826 e le stamperà nel 1829--1830 nelle pubblicazioni dell'università di Kazan', separato dal mondo scientifico, come osserva Eves, da barriere di distanza e di lingua \cite[p. 287]{eves1_1972}. % lang-it
La versione tedesca, \emph{Geometrische Untersuchungen}, la pubblicherà nel 1840. % lang-it
William Kingdon Clifford chiamerà Lobachevskij il ``Copernico della geometria''. % lang-it

% TODO: O geometrii hiperbolicznej wspomina się w $\Delta_{24}^7$. % lang-pl
% TODO: La geometria iperbolica è menzionata in $\Delta_{24}^7$. % lang-it

Eves \cite[s. 291-317]{eves1_1972} rozwija geometrię na płaszczyźnie Łobaczewskiego, gdzie zamiast piątego postulatu przyjmuje się: % lang-pl
Eves \cite[p. 291-317]{eves1_1972} sviluppa la geometria nel piano di Lobachevskij, dove al posto del quinto postulato si assume: % lang-it
\begin{axiom}
[postulat równoległości Łobaczewskiego] % lang-pl
[postulato delle parallele di Lobachevskij] % lang-it
	Niech $P$ będzie punktem poza prostą $AB$, zaś $Q$ jego spodkiem na tę prostą. % lang-pl
	Istnieją wtedy dwie półproste $PX$, $PY$, nieleżące na jednej prostej, nie przecinające prostej $AB$ takie, że dowolna półprosta zawarta w tym kącie $\angle XPY$, w którym leży $PQ$, przecina prostą $AB$. % lang-pl
	Sia $P$ un punto esterno alla retta $AB$ e $Q$ il suo piede su tale retta. % lang-it
	Esistono allora due semirette $PX$, $PY$, non contenute in una stessa retta e che non intersecano la retta $AB$, tali che ogni semiretta contenuta nell'angolo $\angle XPY$ in cui giace $PQ$ interseca la retta $AB$. % lang-it
\end{axiom}

Kąt $\angle XPQ$ nazywamy kątem równoległości w $P$ prostej $AB$. % lang-pl
L'angolo $\angle XPQ$ si chiama angolo di parallelismo in $P$ rispetto alla retta $AB$. % lang-it

Proste przechodzące przez $P$, które nie przechodzą przez ten kąt $\angle XPY$, nazywa hiperrównoległymi. % lang-pl
Istnieje nieskończenie wiele hiperrównoległych do prostej przez punkt poza nią. % lang-pl
Le rette per $P$ che non passano per l'angolo $\angle XPY$ si dicono iperparallele. % lang-it
Esistono infinite rette iperparallele a una retta data per un punto esterno ad essa. % lang-it

Jeśli prosta skierowana $PX$ jest równoległa do prostej skierowanej $AB$, zaś $R$ jest punktem na $PX$ leżącym po tej samej stronie $X$ co $P$, to proste skierowane $RX$ i $AB$ też są równoległe. % lang-pl
Równoległość jest symetryczna: z równoległości $CD$ do $AB$ wynika równoległość $AB$ do $CD$. % lang-pl
Jest też przechodnia: jeśli $AB$ i $CD$ są równoległe do $EF$, to są też równoległe do siebie. % lang-pl
Se la retta orientata $PX$ è parallela alla retta orientata $AB$ e $R$ è un punto di $PX$ che sta dalla stessa parte di $X$ in cui si trova $P$, allora anche le rette orientate $RX$ e $AB$ sono parallele. % lang-it
Il parallelismo è simmetrico: dal parallelismo di $CD$ ad $AB$ segue il parallelismo di $AB$ a $CD$. % lang-it
È anche transitivo: se $AB$ e $CD$ sono parallele a $EF$, allora sono parallele tra loro. % lang-it

Wygodnie jest zdefiniować: % lang-pl
Conviene definire: % lang-it

\begin{definition}
[trójkąt graniczny] % lang-pl
[triangolo limite] % lang-it
	Figurę złożona z dwóch równoległych półprostych oraz odcinka łączącego ich końce nazywamy trójkątem granicznym. % lang-pl
	La figura composta da due semirette parallele e dal segmento che ne congiunge le origini si chiama triangolo limite. % lang-it
\end{definition}

Kąt zewnętrzny trójkąta granicznego jest większy niż przeciwległy kąt wewnętrzny. % lang-pl
Mamy dwie cechy przystawania: bok-kąt (jeśli skończone boki i odpowiednie kąty dwóch trójkątów granicznych są sobie równe, to pozostałe dwie pary kątów też są sobie równe) i kąt-kąt (jeśli dwie pary odpowiednich kątów dwóch trójkątów granicznych są sobie równe, to skończoe boki tych trójkątów też są sobie równe). % lang-pl
L'angolo esterno di un triangolo limite è maggiore dell'angolo interno opposto. % lang-it
Abbiamo due criteri di congruenza: lato-angolo (se i lati finiti e gli angoli corrispondenti di due triangoli limite sono uguali, allora sono uguali anche le altre due coppie di angoli) e angolo-angolo (se due coppie di angoli corrispondenti di due triangoli limite sono uguali, allora sono uguali anche i lati finiti di tali triangoli). % lang-it

Okazuje się, że kąt równoległości w $P$ prostej $AB$ zależy jedynie od odległości $PQ$ i maleje wraz ze wzrostem tej odległości. % lang-pl
Oznaczmy odległość przez $h$, zaś kąt przez $\Pi(h)$. % lang-pl
Przy odpowiednich jednostkach mamy % lang-pl
Si scopre che l'angolo di parallelismo in $P$ rispetto alla retta $AB$ dipende soltanto dalla distanza $PQ$ e diminuisce al crescere di tale distanza. % lang-it
Indichiamo la distanza con $h$ e l'angolo con $\Pi(h)$. % lang-it
Con opportune unità si ha % lang-it
\begin{equation}
	\label{eqn:pi_h}
	\Pi(h) = 2 \arctan \exp (-h).
\end{equation}

W geometrii euklidesowej istnieje naturalna jednostka kątów (na przykład: kąt prosty), ale nie odległości (nie można odtworzyć odcinka jednostkowego po zgubieniu go). % lang-pl
Natomiast w geometrii Łobaczewskiego wzór \ref{eqn:pi_h} pozwala związać odległości z kątami, więc odległości też mają swoją naturalną jednostkę. % lang-pl
Pisze o tym Eves \cite[s. 297]{eves1_1972}. % lang-pl
Nella geometria euclidea esiste un'unità naturale degli angoli (per esempio l'angolo retto), ma non delle distanze (non si può ricostruire il segmento unitario dopo averlo perso). % lang-it
Nella geometria di Lobachevskij, invece, la formula \ref{eqn:pi_h} permette di legare le distanze agli angoli, sicché anche le distanze hanno una loro unità naturale. % lang-it
Ne scrive Eves \cite[p. 297]{eves1_1972}. % lang-it

Eves przechodzi następnie przez czworokąt Saccheriego (\cite[s. 298-301]{eves1_1972}), zależność między polem trójkąta, jego defektem, i równoważnością przez rozcinanie (\cite[s. 301-304]{eves1_1972}). % lang-pl
Defekt trójkąta to różnica między $\pi$ oraz sumą miar jego kątów. % lang-pl
Następnie poruszany jest temat punktów idealnych i~hiperidalnych, odpowiednik punktu w nieskończoności (\cite[s. 304-307]{eves1_1972}; jako zastosowanie tychże, pokazuje, że symetralne boków trójkąta są współpękowe). % lang-pl
Eves passa poi al quadrilatero di Saccheri (\cite[p. 298-301]{eves1_1972}), al legame tra l'area di un triangolo, il suo difetto e l'equiscomposizione (\cite[p. 301-304]{eves1_1972}). % lang-it
Il difetto di un triangolo è la differenza tra $\pi$ e la somma delle misure dei suoi angoli. % lang-it
Segue il tema dei punti ideali e ultraideali, analoghi del punto all'infinito (\cite[p. 304-307]{eves1_1972}; come applicazione di essi dimostra che gli assi dei lati di un triangolo sono concorrenti). % lang-it

W sekcji ósmej Eves przytoczy twierdzenie Hjemsleva (czyli nasze twierdzenie \ref{thm:hjelmslev}) i wprowadzi specjalne przekształcenie płaszczyzny: % lang-pl
Nella sezione ottava Eves riporterà il teorema di Hjelmslev (cioè il nostro teorema \ref{thm:hjelmslev}) e introdurrà una speciale trasformazione del piano: % lang-it

\begin{definition}
[przekształcenie Hjelmsleva] % lang-pl
[trasformazione di Hjelmslev] % lang-it
	Niech $O$ będzie ustalonym punktem, zaś $\theta$ ustalonym ostrym kątem skierowanym. % lang-pl
	Przekształcenie $S$ przenosi punkt $P$ na punkt $P'$ taki, że $\angle POP' = \theta$ oraz $\angle OP'P = \pi/2$. % lang-pl
	Złożenie obrotu wokół $O$ o kąt $-\theta$ z $S$ nazywamy przekształceniem Hjelmsleva i oznaczamy $T$. % lang-pl
	Sia $O$ un punto fissato e $\theta$ un angolo acuto orientato fissato. % lang-it
	La trasformazione $S$ manda il punto $P$ nel punto $P'$ tale che $\angle POP' = \theta$ e $\angle OP'P = \pi/2$. % lang-it
	La composizione di una rotazione intorno a $O$ di angolo $-\theta$ con $S$ si chiama trasformazione di Hjelmslev e si indica con $T$. % lang-it
\end{definition}

\begin{proposition}
	Przekształcenie Hjelmsleva na płaszczyźnie euklidesowej jest jednokładnością $H(O, \cos \theta)$. % lang-pl
	Nel piano euclideo la trasformazione di Hjelmslev è l'omotetia $H(O, \cos \theta)$. % lang-it
\end{proposition}

\begin{proposition}
	Przekształcenie Hjelmsleva przenosi płaszczyznę Łobaczewskiego na wnętrze okręgu o środku $O$ i promieniu $h$, gdzie $\Pi(h) = \theta$. % lang-pl
	La trasformazione di Hjelmslev manda il piano di Lobachevskij nell'interno del cerchio di centro $O$ e raggio $h$, dove $\Pi(h) = \theta$. % lang-it
\end{proposition}

Przekształcenie $T$ pozwala uzasadnić wiele właściwości płaszczyzny Łobaczewskiego, ale z braku miejsca, Eves odstąpi od tego. % lang-pl
La trasformazione $T$ permette di giustificare molte proprietà del piano di Lobachevskij, ma per mancanza di spazio Eves vi rinuncerà. % lang-it

%